
\section{Introduction}

Since before the COVID-19 pandemic, various CS and non-CS faculty at
our North American postsecondary higher-education
institution had been promoting mastery 
learning as a way of both improving 
pedagogy~\cite{Bloom68} and narrowing equity gaps~\cite{guskey2007closing}.
Since trustworthy and efficient summative assessment is an important
component of mastery learning,
we aimed to pilot a Computer-Based Testing
Facility (CBTF) modeled on the successful
experiences at
the University of Illinois at Urbana-Champaign~\cite{zilles-testing-less-trying} and
other institutions~\cite{downey-cbtf}.
Our goal was to
start with a single course and eventually spread the benefits of
mastery learning to all departments throughout the institution.
Because we believed that such a testing center should be a common good
and therefore centrally administered and resourced, scale-up would
require demonstrating the utility of the service beyond
engineering/STEM courses and departments,
centralizing its operations in a nonacademic campus unit,
and developing cross-unit funding and coordination models that did not
currently exist.

Our ``proto-CBTF'' pilot in Fall 2023 consisted of a single course with
230 students consuming about 1,000 student-exam-hours and proctored
largely by TAs and one-off student assistants.
2.5 years later, we have successfully scaled to 24 courses from 16 departments, enrolling over
5,000 undergraduates and accounting for over 45,000
seated exam-hours, proctored by a staff of 29 regular proctors (also student assistants).
These students take exams in one dedicated 30-seat CBTF and a second
shared 30-seat lab; the campus has since committed to opening a
dedicated 80--100 seat CBTF in Fall 2026 in a renovated building, 
and that new CBTF is already fully subscribed, months before opening.
All of this is now offered as  a ``common good'' campus service operated
by a comprehensive and
well-managed instructional technology unit, which also
has mature processes for hiring, onboarding, and supervision of
proctors (student assistants).
This paper describes the ``proto-CBTF'' bootstrap period with a few
details about the ``scale up'' period that led to the present moment.

The single course chosen for our pilot was an upper-division
software/systems course enrolling around 230 students.
During the COVID-19 pandemic, the course's quizzes had been converted
to the PrairieLearn assessment authoring
system\footnote{prairielearn.org} and were taken remotely using
videoconference-based proctoring.  Once the pandemic lockdown ended,
we decided to convert the 2.5-hour final exam (traditionally a written
final with a mix of short-answer and long-form questions) to be administered
in a prototype CBTF in Fall 2023.

A key property of the ``UIUC model'' we were trying to emulate
is that the exam incorporates enough randomness~\cite{chen-how-much-randomization} that its
integrity does not depend heavily on all
students take it simultaneously---a necessity since we had only about
57 seats in our proto-CBTF spaces, and a boon to students since they could
self-schedule their exam over a multi-day window.
The model also embodies other
practices that are important at scale but could  have been ignored for our single course;
since we planned to scale, we tried to avoid scale-thwarting
practices, indicating in the text
where we would have done things differently in hindsight.

We divide our experience into three categories: physical space and
configuration (Section~\ref{sec:physical}), IT infrastructure
(Section~\ref{sec:it}), and policies/coordination/communication (Section~\ref{sec:policy}).
Due to space constraints and our focus on operational lessons, we do
not discuss exam content.  In particular, we do not discuss the
process of moving from our paper-based exam to PrairieLearn, except to
state that we ended up with an exam that was in many ways as rich and
authentic as our paper-based final, and in some ways significantly
better.

